\subsection*{4.3 Lower bounds}
In this section we show that if one tries to get a program of length less than $3n$, one cannot beat Theorem3.
\begin{proof}[Proof of Theorem4]
Let $P=\pi_1\pi_2\ldots\pi_s$ be any program computing $\mathrm{AND}(x_1,\ldots,x_n)$. We will write $\sigma_i^j$ to refer to the permutation in the $j$th instruction in $P$ that reads variable $x_i$; in other words, the instructions corresponding to $x_i$ will be $\langle i,\sigma_i^1\rangle\ldots\langle i,\sigma_i^k\rangle$ for some $k$.
\paragraph{Claim27.} Any program $P$ computing AND of more than one variable must read every variable at least twice.
\begin{proof}
Assume that some variable $x_i$ is read only once in $P$. Then setting $x_{i'}=0$ for all $i'\neq i$, we get $\sigma_i^1=P(0,\ldots,0,1,0,\ldots,0)=\mathrm{id}$. Therefore $P$ acts identically whether $x_i$ is0 or1, which is a contradiction because AND depends on $x_1$.
\end{proof}
Now consider when some variable $x_i$ is read exactly twice. Then $P$ has the form:
\[\Sigma_1,\langle i,\sigma_i^1\rangle,\Pi_1,\langle i,\sigma_i^2\rangle,\Sigma_2.\]
Rotate this program to produce
\[P'=\langle i,\sigma_i^1\rangle,\Pi_1,\langle i,\sigma_i^2\rangle,\Pi_2\]
where $\Pi_2=\Sigma_2,\Sigma_1$. By Lemma22, $P'$ also computes AND.
The following is our main claim.
\paragraph{Claim28.} Every variable besides $x_i$ is read at least once in $\Pi_1$, and there is at most one such variable $x_{i'}$ which is not read at least twice in $\Pi_1$. The same holds for $\Pi_2$.
\begin{proof}
First, assume for contradiction that there exists $i'\neq i$ such that $x_{i'}$ does not appear in $\Pi_1$. Then if we fix $x_{i''}=1$ for all $i''\neq i,i'$, we can apply Corollary21 to move all instructions querying variables other than $x_i$ and $x_{i'}$ to the end of the program, to get an equivalent program of the following form which computes $\mathrm{AND}(x_i,x_{i'})$:
\[\langle i,\sigma_i^{\prime1}\rangle,\langle i,\sigma_i^{\prime2}\rangle,\Sigma\]
where $\Sigma$ is a sequence of instructions which do not query $x_i$. Applying Lemma18, we can replace $\langle i,\sigma_i^{\prime1}\rangle,\langle i,\sigma_i^{\prime2}\rangle$ with a single instruction $\langle i,\sigma''\rangle$ to we get an equivalent program
\[\langle i,\sigma''_i\rangle,\Sigma\]
which also computes $\mathrm{AND}(x_i,x_{i'})$, which contradicts Claim27 as $i$ is only read once.
Next, assume for contradiction that there exist $i'\neq i''\neq i$ such that $i'$ and $i''$ appear only once each in $\Pi_1$. If we fix $x_{i'''}=1$ for all $i'''\neq i,i',i''$, applying Lemmas21 and18, without loss of generality the following program computes $\mathrm{AND}(x_i,x_{i'})$:
\[\langle i,\sigma_i^{\prime1}\rangle,\langle i',\sigma_{i'}\rangle,\langle i'',\sigma_{i''}\rangle,\langle i,\sigma_i^{\prime2}\rangle,\Sigma\]
where $\sigma_{i'}$ and $\sigma_{i''}$ are some permutations and $\Sigma$ is a set of instructions reading only the variables $x_{i'}$ and $x_{i''}$.
Define $\Sigma_{i'}$ to be the result of fixing $x_{i''}=0$ in $\Sigma$, and define $\Sigma_{i''}$ to be the result of fixing $x_{i'}=0$ in $\Sigma$. Note that if only one remaining variable is set to1 then the program must output0, so $\sigma_i^{\prime2}=(\sigma_i^{\prime1})^{-1}$, $\Sigma_{i'}=(\sigma_{i'})^{-1}$, and $\Sigma_{i''}=(\sigma_{i''})^{-1}$. Thus if we set $x_{i''}=0$ our resulting program is
\[\langle i,\sigma_i^{\prime1}\rangle,\langle i',\sigma_{i'}\rangle,\langle i,(\sigma_i^{\prime1})^{-1}\rangle,\langle i',(\sigma_{i'}^1)^{-1}\rangle\]
and so setting $x_i=x_{i'}=1$ we get that $\sigma_i^{\prime1}\sigma_{i'}(\sigma_i^{\prime1})^{-1}(\sigma_{i'})^{-1}=\mathrm{id}$. Therefore by Lemma19 we can swap the order of these two instructions and get an equivalent program for $\mathrm{AND}(x_i,x_{i'},x_{i''})$ of the form
\[\langle i',\sigma_{i'}\rangle,\langle i,\sigma_i^{\prime1}\rangle,\langle i'',\sigma_{i''}\rangle,\langle i,(\sigma_i^{\prime1})^{-1}\rangle,\Sigma\]
and similarly by fixing $x_{i'}=0$ we have $\sigma_i^{\prime1}\sigma_{i''}(\sigma_i^{\prime1})^{-1}(\sigma_{i''})^{-1}=\mathrm{id}$, which by Lemma19 leaves us with the program
\[\langle i',\sigma_{i'}\rangle,\langle i'',\sigma_{i''}\rangle,\langle i,\sigma_i^{\prime1}\rangle,\langle i,(\sigma_i^{\prime1})^{-1}\rangle,\Sigma\]
and applying Lemma18 on our two layers reading $i$ gives us a program which never reads $x_i$, which is a contradiction.
Finally, to prove the claim for $\Pi_2$, observe that by Lemma22,
\[P''=\langle i,\sigma_i^2\rangle,\Pi_2,\langle i,\sigma_i^1\rangle,\Pi_1\]
also computes AND, and apply the above argument to $P''$, with $\Pi_2$ playing the role of $\Pi_1$.
\end{proof}
By a simple analysis of Claim28, one of two cases must occur for the variables besides $x_i$: either1) one variable $x_{i'}$ is read at least twice and all other variables are read at least four times or2) two variables $x_{i'},x_{i''}$ are read at least three times and all other variables are read at least four times. This is because a read-2 variable can only be read at most once in each of $\Pi_1$ and $\Pi_2$, while a read-3 variable will be read at most once in either $\Pi_1$ or $\Pi_2$. In either of these cases, our branching program must have length at least $4(n-2)+2\cdot2=4(n-3)+2\cdot1+3\cdot2=4n-4$.
\end{proof}
